\chapter{Conclusioni}
\label{cap:conclusioni}

Il presente elaborato ha avuto come obiettivo principale l'analisi e la valutazione di soluzioni per la generazione di report, con particolare attenzione alla possibilità di realizzare un sistema personalizzato basato sull'utilizzo del motore BIRT. A partire da un'analisi preliminare delle principali piattaforme di reporting, è stato progettato e sviluppato un Proof of Concept (PoC) volto a dimostrare la fattibilità di un approccio alternativo rispetto a soluzioni consolidate come JasperServer.\\
Il lavoro svolto ha previsto diverse fasi: una prima fase di studio delle tecnologie esistenti, una fase di progettazione dell'architettura software e infine una fase di implementazione e testing. In particolare, è stata sviluppata la componente centrale \texttt{BirtDesignToDocument}, responsabile della generazione dei documenti a partire da file di design BIRT e dati in formato JSON, integrata all'interno di un'architettura di tipo Model-View-Controller.

\section{Risultati ottenuti}

I risultati ottenuti dimostrano che il Proof of Concept sviluppato è in grado di soddisfare i requisiti funzionali definiti nella fase di analisi. In particolare, il sistema consente di generare report in diversi formati (PDF, DOC, XLSX, HTML) a partire da dati strutturati, garantendo flessibilità e indipendenza da piattaforme server esterne.\\
Dal punto di vista prestazionale, i test condotti hanno evidenziato come il comportamento del sistema sia complessivamente comparabile a quello di JasperServer, soprattutto nelle esecuzioni successive alla prima. In entrambi i casi, infatti, il tempo di inizializzazione rappresenta un fattore significativo, mentre le esecuzioni successive risultano più efficienti grazie alla riduzione dell'overhead iniziale.\\
Un ulteriore risultato rilevante riguarda la modularità del sistema sviluppato: la separazione tra logica di generazione, configurazione e interfaccia consente una maggiore flessibilità e facilita eventuali estensioni future. Questo aspetto rappresenta un valore aggiunto rispetto a soluzioni più monolitiche.

\section{Limiti del lavoro}

Nonostante i risultati positivi, il lavoro presenta alcuni limiti che è importante evidenziare.
In primo luogo, i test prestazionali sono stati condotti in un ambiente controllato e non rappresentano pienamente scenari di produzione reale. In particolare, non sono stati eseguiti test sotto carico concorrente, che potrebbero evidenziare criticità legate alla gestione delle risorse e alla scalabilità del sistema.\\
In secondo luogo, il numero di scenari testati è limitato e non copre tutte le possibili casistiche applicative. Ad esempio, non sono stati considerati dataset di dimensioni molto elevate o configurazioni distribuite.\\
Un ulteriore limite riguarda la gestione delle risorse del motore BIRT: l'inizializzazione e la chiusura del motore ad ogni esecuzione, seppur corretta dal punto di vista della stabilità, potrebbe risultare inefficiente in contesti ad alta frequenza di utilizzo.\\
Infine, l'interfaccia utente sviluppata nel PoC ha uno scopo puramente dimostrativo e non è stata progettata per un utilizzo in ambienti produttivi.

\section{Sviluppi futuri}

Alla luce dei risultati ottenuti, è possibile individuare diverse direzioni per sviluppi futuri del sistema.\\
Un primo possibile miglioramento consiste nell'ottimizzazione delle prestazioni, ad esempio attraverso l'introduzione di meccanismi di caching o la gestione persistente del motore BIRT, evitando la sua inizializzazione ad ogni esecuzione.\\
Un secondo sviluppo riguarda l'evoluzione dell'architettura verso un modello a servizi, ad esempio tramite la realizzazione di un'API REST che consenta di integrare il sistema in applicazioni web o microservizi.\\

Ulteriori miglioramenti potrebbero includere:
\begin{itemize}
    \item supporto a nuove sorgenti dati (database relazionali, servizi esterni);
    \item gestione della concorrenza e ottimizzazione per scenari multi-utente;
    \item miglioramento dell'interfaccia utente;
    \item integrazione con sistemi aziendali esistenti.
\end{itemize}

Questi sviluppi consentirebbero di trasformare il Proof of Concept in una soluzione completa e utilizzabile in contesti reali.

\section{Raggiungimento degli obiettivi}
Di seguito si riporta la verifica dei requisiti individuati durante l'analisi descritta nel Capitolo \ref{cap:analisi-requisiti}, suddivisi in funzionali, non funzionali, qualitativi e di vincolo. Per ciascun requisito è indicato l'esito finale ottenuto al termine del progetto: soddisfatto, parzialmente soddisfatto oppure non soddisfatto.

\subsection{Requisiti funzionali}
\begin{table}[H]%
            \caption{Tabella del tracciamento dei requisti funzionali}
            \label{tab:requisiti-funzionali-conclusioni}
            \begin{tabularx}{\textwidth}{|lXll|}
            \hline\hline
            \rowcolor{gray!40}
            \textbf{Requisito} & \textbf{Descrizione} & \textbf{Classificazione} & \textbf{Esito}\\
            \hline
            RFO1    & Il sistema deve fornire un editor grafico per consentire la creazione e la progettazione visuale dei report & Obbligatorio & Soddisfatto\\
            \hline
            RFO2    & Il sistema deve permettere l'integrazione con servizi e applicazioni che espongono dati tramite \emph{\gls{apig}} in formato \emph{JSON} & Obbligatorio & Soddisfatto\\
            \hline
            RFO3    & Il sistema deve poter gestire report complessi e strutturati in modo modulare, in particolare elementi di sottoreport & Obbligatorio & Soddisfatto\\
            \hline
            RFO4    & Il sistema deve consentire l'esportazione dei report in molteplici formati (PDF, DOC, XLSX) & Obbligatorio & Soddisfatto\\
            \hline
            RFO5    & Il sistema deve consentire l'invocazione programmatica dei report da applicazioni esterne tramite \emph{\gls{apig}} & Obbligatorio & Soddisfatto\\
            \hline
            RFO6    & Il sistema deve garantire supporto a database \emph{SQL} e \emph{NoSQL} & Obbligatorio & Soddisfatto\\
            \hline
            RFO7    & Il sistema deve fornire la possibilità di personalizzazione e \emph{scripting} dei report & Obbligatorio & Soddisfatto\\
            \hline
            \end{tabularx}
        \end{table}%

\subsection{Requisiti non funzionali}
\begin{table}[H]%
            \caption{Tabella del tracciamento dei requisiti non funzionali}
            \label{tab:requisiti-non-funzionali-conclusioni}
            \begin{tabularx}{\textwidth}{|lXll|}
            \hline\hline
            \rowcolor{gray!40}
            \textbf{Requisito} & \textbf{Descrizione} & \textbf{Classificazione} & \textbf{Esito}\\
            \hline
            RNFO1    & Il sistema deve supportare la scalabilità, consentendo la gestione di un numero crescente di richieste e utenti & Obbligatorio & Soddisfatto\\
            \hline
            RNFO2    & Il sistema deve garantire la corretta generazione dei report in modo consistente & Obbligatorio & Soddisfatto\\
            \hline
            RNFO3    & Il sistema deve consentire una facile integrazione con applicazioni esistenti sviluppate in \emph{Java}, \emph{.NET} e \emph{Python} & Obbligatorio & Soddisfatto\\
            \hline
            RNFO4    & Il sistema deve garantire tempi di risposta accettabili anche su dataset consistenti & Obbligatorio & Soddisfatto\\
            \hline
            RNFO5    & Il sistema deve risultare mantenibile, facilitando aggiornamenti, configurazioni e interventi evolutivi & Desiderabile & Soddisfatto\\
            \hline
            \end{tabularx}
        \end{table}%
        
\subsection{Requisiti qualitativi}
\begin{table}[H]%
            \caption{Tabella del tracciamento dei requisiti qualitativi}
            \label{tab:requisiti-qualitativi-conclusioni}
            \begin{tabularx}{\textwidth}{|lXll|}
            \hline\hline
            \rowcolor{gray!40}
            \textbf{Requisito} & \textbf{Descrizione} & \textbf{Classificazione} & \textbf{Esito}\\
            \hline
            RQO1    & Il sistema deve fornire un'interfaccia utente intuitiva e di facile utilizzo & Obbligatorio & Soddisfatto\\
            \hline
            RQD2    & Il sistema deve garantire coerenza nella progettazione e nel \emph{layout} dei report & Desiderabile & Soddisfatto\\
            \hline
            RQD3    & Il sistema deve offrire una buona esperienza di sviluppo, supportata da documentazione chiara e \emph{\gls{apig}} ben definite & Desiderabile & Soddisfatto\\
            \hline
            \end{tabularx}
        \end{table}%

\subsection{Requisiti di vincolo}
\begin{table}[H]%
            \caption{Tabella del tracciamento dei requisiti di vincolo}
            \label{tab:requisiti-vincolo-conclusioni}
            \begin{tabularx}{\textwidth}{|lXll|}
            \hline\hline
            \rowcolor{gray!40}
            \textbf{Requisito} & \textbf{Descrizione} & \textbf{Classificazione} & \textbf{Esito}\\
            \hline
            RVO1    & Il sistema deve poter essere installato e gestito su server locali, senza dipendere esclusivamente da soluzioni cloud & Obbligatorio & Soddisfatto\\
            \hline
            RVO2    & Il sistema deve prevedere condizioni di \emph{licensing} sostenibili, evitando vincoli restrittivi legati a licenze proprietarie & Obbligatorio & Soddisfatto\\
            \hline
            RVO3    & Il sistema deve essere compatibile con l'infrastruttura e lo \emph{stack} tecnologico aziendale esistente & Desiderabile & Soddisfatto\\
            \hline
            RVD4    & Il sistema dovrebbe utilizzare tecnologie consolidate e ampiamente supportate dalla \emph{community} o dal \emph{vendor} & Desiderabile & Soddisfatto\\
            \hline
            \end{tabularx}
        \end{table}%
\subsection{Riepilogo dei requisiti raggiunti}
\begin{table}[H]
            \centering
            \caption{Riepilogo dei requisiti soddisfatti}
            \label{tab:riepilogo-requisiti-conclusioni}
            \begin{tabular}{|l|c|c|c|c|}
            \hline\hline
            \rowcolor{gray!40}
            \textbf{Tipologia} & \textbf{Obbligatorio} & \textbf{Desiderabile} & \textbf{Opzionale} & \textbf{Totale}\\
            \hline
            \textbf{Funzionali} & 7 & 0 & 0 & 7\\
            \hline
            \textbf{Non funzionali} & 4 & 1 & 0 & 5\\
            \hline
            \textbf{Qualitativi} & 1 & 2 & 0 & 3\\
            \hline
            \textbf{Vincolo} & 3 & 1 & 0 & 4\\
            \hline
            \textbf{Totale} & 15 & 4 & 0 & 19\\
            \hline
            \end{tabular}
        \end{table}
Dalla tabella illustrata si può notare che tutti i requisiti identificati durante l'analisi sono stati soddisfatti, con un totale di 19 requisiti raggiunti. In particolare, sono stati soddisfatti tutti i requisiti funzionali,non funzionali, qualitativi e di vincolo. Questo risultato conferma il successo del progetto nel raggiungere gli obiettivi prefissati e nel fornire una soluzione efficace per la generazione di report basata su BIRT.

\section{Considerazioni finali}

In conclusione, il lavoro svolto ha dimostrato la fattibilità della realizzazione di un sistema di reportistica basato su BIRT completamente personalizzato e indipendente da piattaforme server esterne.\\
Il Proof of Concept sviluppato rappresenta una valida alternativa a soluzioni consolidate, almeno per quanto riguarda scenari di utilizzo standard, offrendo al contempo maggiore flessibilità e controllo sull'intero processo di generazione dei report.\\
L'esperienza maturata durante lo sviluppo ha inoltre permesso di approfondire tematiche rilevanti nell'ambito dell'ingegneria del software, quali la progettazione architetturale, l'integrazione tra componenti e l'analisi delle prestazioni.\\
Nel complesso, i risultati ottenuti confermano il raggiungimento degli obiettivi prefissati e forniscono una base solida per eventuali evoluzioni future del sistema.

\section{Valutazione personale}
Il progetto di stage ha rappresentato un'importante opportunità di crescita professionale e personale, consentendomi di applicare le conoscenze acquisite durante il percorso universitario a un contesto reale e complesso. \\
Ringrazio l'azienda per avermi dato la possibilità di lavorare su un progetto così stimolante e per il supporto ricevuto durante tutto il periodo di stage.\\
Ho avuto modo di confrontarmi con sfide tecniche significative, sviluppando competenze specifiche nell'ambito della generazione di report e dell'integrazione di sistemi software, nonché migliorando le mie capacità di problem solving, gestione del tempo e lavoro in team. \\
Con questa esperienza poi ho potuto migliorare le mie capacità di analisi, progettazione e sviluppo software, permettendomi di riuscire a conciliare attività di ricerca e sviluppo con le esigenze di un progetto concreto, rispettando scadenze e obiettivi prefissati.\\
Inoltre, l'esperienza mi ha permesso di capire meglio quale anmbito dell'informatica sia più adatto alle mie competenze e alla mia persona. \\
Ci sono statti anche dei limiti e delle difficoltà, specie all'inizio del progetto, quando ho dovuto affrontare una curva di apprendimento piuttosto ripida per comprendere le tecnologie e gli strumenti utilizzati. Tuttavia, grazie al supporto del team e alla mia determinazione, sono riuscito a superare queste difficoltà e a portare a termine con successo il progetto.\\
Nel complesso sono soddisfatto dei risultati ottenuti e grato per l'opportunità di aver potuto contribuire a questo progetto di ricerca.